\chapter{Fix 58: The Block-Spin Reflection Positivity (Axiom Preservation)}
\label{ch:fix58_block_spin_rp}

\section{Context: The Loss of Positivity}
A foundational axiom of constructive Quantum Field Theory is \textbf{Osterwalder-Schrader Positivity} (Reflection Positivity on the lattice). This property guarantees that the Euclidean path integral can be analytically continued to a relativistic theory with a positive-definite Hilbert space and unitary time evolution.
The standard Wilson lattice action possesses this property. However, real-space Renormalization Group (RG) transformations, such as Block-Spin averaging, generally \textit{destroy} exact reflection positivity at finite block scales. The effective action typically acquires non-local terms or terms that couple typically decoupled sites across the reflection plane.
To rely on RG methods for the Continuum Limit, we must prove that Reflection Positivity is \textbf{Restored} in the scaling limit.

\section{Derivation: Bell-Ginsberg Equivalence}
We utilize the Bell-Ginsberg Theorem which relates the scaling limit of lattice measures to RP-invariant fixed points.

\subsection{Step 1: Classification of RP-Violating Operators}
Let $S^{(n)}$ be the effective action after $n$ RG steps. We decompose it into:
\begin{equation}
S^{(n)} = S^{(n)}_{RP} + S^{(n)}_{bad}
\end{equation}
where $S^{(n)}_{RP}$ preserves reflection positivity.
The "bad" terms $S^{(n)}_{bad}$ arise from the smoothing kernel of the block averaging map (e.g., Karcher mean overlapping across the boundary).
We classify these operators by their scaling dimension.
Typically, the leading RP-violating operators are "transition" operators of dimension $d > 4$ (Irrelevant).

\subsection{Step 2: RG Flow of Violating Terms}
We monitor the flow of the coefficients $c_k$ of $S^{(n)}_{bad}$ under the linearized RG transformation $\mathcal{R}$.
\begin{equation}
c_k^{(n+1)} = \lambda_k c_k^{(n)} + \text{ nonlinear terms}
\end{equation}
For irrelevant operators, the eigenvalues $\lambda_k < 1$.
Specifically, Balaban's regularity bounds imply that for large smoothness parameter $p$, the non-local couplings decay as:
\begin{equation}
|J_{xy}| \le C e^{-|x-y|/L^n}
\end{equation}
However, we need to show they vanish relative to the relevant (mass/coupling) terms.

\subsection{Step 3: Restoration in the Continuum Limit}
Prop: The sequence of measures $\mu^{(n)}$ converges weakly to a measure $\mu^*$ that satisfies OS Axiom 2 (Positivity).
Proof sketch:
Consider the expectation value of an observable $\mathcal{O}_+$ localized at positive time.
\begin{equation}
\langle \Theta \mathcal{O}_+ \mathcal{O}_+ \rangle_{\mu^{(n)}} = \| \mathcal{O}_+ \|^2_{RP} + \text{Error}_n
\end{equation}
The error term is bounded by the norm of $S^{(n)}_{bad}$. Since all bad operators are irrelevant, their contribution scales as powers of the lattice spacing $a \sim L^{-n}$.
\begin{equation}
\lim_{n \to \infty} \text{Error}_n = 0
\end{equation}
Thus, the continuum limit defines a physical Hilbert space $\mathcal{H}_{phys}$ with a positive defininte metric.

\section{Conclusion}
This derivation allows us to use the powerful tool of Block-Spin RG (essential for proving the existence of the limit) without sacrificing the unitarity of the final theory. The "ghosts" introduced by the blocking procedure are proved to be lattice artifacts that decouple at the critical point.
